% ======================================================================
%  CAPA — UFBA / Escola Politécnica / Eng. Química / PRH
% ======================================================================

% --- Fonte da capa ---
\newcommand{\capaFamilia}{\rmfamily} % Força a fonte Palatino do seu texto

% --- Cores das faixas laterais --- 
\definecolor{capaAzul}{HTML}{4472C4}
\definecolor{capaVerde}{HTML}{26F22B}

% --- Geometria das faixas (medidas a partir da borda esquerda da folha) ---
\newcommand{\faixaAzulInicio}{18.7mm}
\newcommand{\faixaAzulFim}{29.0mm}
\newcommand{\faixaVerdeInicio}{28.8mm}
\newcommand{\faixaVerdeFim}{33.2mm}

% --- Caminhos dos logos (Ajustados para a sua pasta de imagens) ---
\newcommand{\logoPRH}{imagens/logo_prh_pequeno}
\newcommand{\logoUFBA}{imagens/logo_ufba}
\newcommand{\logoANP}{imagens/logo_anp}
\newcommand{\logoPRHprog}{imagens/logo_prh_circulo}
\newcommand{\logoFinep}{imagens/logo_finep}

\newcommand{\logo}[3]{%
  \IfFileExists{#1.svg}{\includegraphics[#2]{#1}}{%
  \IfFileExists{#1.pdf}{\includegraphics[#2]{#1}}{%
  \IfFileExists{#1.png}{\includegraphics[#2]{#1}}{%
  \IfFileExists{#1.jpg}{\includegraphics[#2]{#1}}{%
    \fbox{\parbox[c][2.2cm][c]{2.8cm}{\centering\footnotesize\ttfamily #3}}%
  }}}}%
}

\thispagestyle{empty}

% Faixas verticais à esquerda e posicionamento absoluto do logo PRH
\begin{tikzpicture}[remember picture, overlay]
  % Faixa Azul
  \fill[capaAzul]
    ([xshift=\faixaAzulInicio]current page.north west) rectangle
    ([xshift=\faixaAzulFim]current page.south west);
  % Faixa Verde
  \fill[capaVerde]
    ([xshift=\faixaVerdeInicio]current page.north west) rectangle
    ([xshift=\faixaVerdeFim]current page.south west);
    
  % SOLUÇÃO: Posiciona o logo do PRH exatamente no centro geométrico das listras (25.95mm)
  % O yshift=-1.4cm alinha o topo dele perfeitamente com o topo dos outros logos da direita.
  \node[anchor=north, inner sep=0pt] at ([xshift=31.95mm, yshift=-2.0cm]current page.north west) {
    \logo{\logoPRHprog}{height=4.5cm}{PRH 36}
  };
\end{tikzpicture}

% 1. Puxa o início da área útil para cima (combina com o yshift da listra)
\vspace*{-1.6cm}

% 2. Afasta todo o bloco de texto para a direita (fora das listras) e limita a largura
\hspace*{1.4cm}%
\begin{minipage}{\dimexpr\textwidth-3cm\relax}
\capaFamilia
\centering

% ---------- LOGOS DA DIREITA ----------
% Como o PRH foi para a listra, os outros 4 logos ficam centralizados na área branca
\vspace*{0cm} 
{\centering
  \renewcommand{\arraystretch}{1.6}%
  \begin{tabular}{cc}
    \logo{\logoUFBA}{height=2.2cm}{UFBA} &
    \logo{\logoANP}{height=1.6cm}{anp} \\
    \logo{\logoPRH}{height=1.4cm}{prh} &
    \logo{\logoFinep}{height=1.4cm}{Finep} \\
  \end{tabular}\par
}%

% ---------- CABEÇALHO INSTITUCIONAL ----------
\vspace{0.9cm}
{\Large\bfseries
    UNIVERSIDADE FEDERAL DA BAHIA\\
    ESCOLA POLITÉCNICA\\
    CURSO DE GRADUAÇÃO EM ENGENHARIA QUÍMICA}

\vspace{0.2cm}
{\large\bfseries
  PROGRAMA DE RECURSOS HUMANOS EM PETRÓLEO E MEIO AMBIENTE DA UNIVERSIDADE FEDERAL DA BAHIA\par}

% ---------- AUTOR ----------
\vspace{1.5cm}
{\large\bfseries \imprimirautor \par}

% ---------- TÍTULO ----------
\vspace{1.5cm}
{\large\bfseries
  \imprimirtitulo \par}

% ---------- ORIENTADOR ----------
\vspace{1.5cm}
{\large \imprimirorientador \par}

% ---------- LOCAL E ANO ----------
\vspace{0.5cm}
{\large\bfseries
  SALVADOR-BA\\
  2026\par}

\end{minipage}